\documentclass[11pt,a4paper]{article}
\usepackage[margin=22mm,top=19mm,bottom=20mm]{geometry}
\usepackage[T1]{fontenc}
\usepackage{lmodern,microtype,amsmath}
\usepackage[hidelinks,pdftitle={Move32}]{hyperref}
\pagestyle{plain}
\setlength{\parindent}{0pt}
\setlength{\parskip}{7pt}
\emergencystretch=1em

\begin{document}
{\LARGE\bfseries Move32}\par
{\small Working design / 1 November 2026 / companion to Rule42}

\textbf{Shared contribution}\par
\emph{Risk is a budget, not a penalty.} A person states how often failure is
acceptable. Evolution ranks candidates by whether they meet that budget, then by
how well they do the task. No safety weight is tuned, no cost function is designed
beyond the violation itself.

\bigskip
{\Large\bfseries The idea, and the two studies that carry it}\par
{\small Page 1: framing}

\textbf{What the brief asks.} The allocated topic asks whether evolutionary
algorithms can learn safe policies without hand-crafted penalties or cost
functions, balancing reward against the probability of violation. Penalty methods
need a weight sweep: too weak is unsafe, too strong is useless. Lagrangian
reinforcement learning learns the weight but needs many samples. Our answer is to
replace the weight with a number a human can actually state.

\textbf{The rule.} For a candidate adapter with parameters $\theta$, evaluated on
$n$ complete episodes, record the mean task return $\widehat J(\theta)$ and the
fraction of episodes with any violation $\widehat p_V(\theta)$. At a requested
limit $\epsilon$, candidates with $\widehat p_V\le\epsilon$ outrank all others;
among them, higher $\widehat J$ wins; among the rest, smaller shortfall wins.
This is the fitness. Everything else is a control.

\textbf{The dial.} The same rule at several requested limits gives a library of
controllers. The claim is tested by one figure: requested risk on one axis,
measured risk and completion on the other, one curve per method. Stricter
requests should give lower measured risk with useful completion retained. We do
not assume monotonicity and we do not pick limits after seeing the curve.

\textbf{What evolution contributes.} Not optimizer superiority. Our data say random
search ties CMA-ES in small adapter families. The contribution is constraint
handling under noise: the constraint is a probability estimated from a handful of
episodes, and we ask whether feasibility-first ranking lands near the requested
risk without a weight sweep, against penalty-weighted evolution and Lagrangian RL
at matched simulator budgets.

\textbf{The vehicle.} A frozen pretrained system plus a compact adapter with
generic initialisation and no expert fallback. The frozen system makes episodes
cheap and the base behaviour competent. It is the setting, not the claim.

\textbf{LeWM's job.} Develop the recipe fast. Episodes cost seconds, so the
ranking rule, the dial and the matched comparison are debugged there, on TwoRoom
with Anchored Planning as the base planner. PushT and Cube are thesis material if
they stay cheap.

\textbf{Block Plucking's job.} Show the recipe means something physical. A frozen
robot policy that drops blocks. A real failure predicate, not a drawn circle. Two
hand-tuned references already published: the UNISafe safety filter and the task 9
release rule. The concrete question: can a small evolved adapter, given only a
risk budget, reach the completion of those hand-designed fixes at the risk level
we asked for?

\textbf{One experiment at two scales.} Nothing in Block Plucking is new method.
Same adapter family, same ranking rule, same arms, same figure. Only the base
system and the failure are real.

\newpage
{\Large\bfseries Experiment 2: a risk-budgeted adapter around frozen UNISafe}\par
{\small Page 2: design}

\textbf{The frozen system, as it exists.} Franka arm, three-block stack, extract
the middle block. 300 controls at 20 Hz, seven-coordinate actions, six arm deltas
and one gripper sign. Dreamer world model, task actor, uncertainty ensemble and
reachability filter, all frozen, on the migrated Isaac Sim 6.2 stack validated on
25 October. Native completes about 3/64 reserved starts and violates in about
6/64. Hand-designed release rules complete 54 to 56 of 64 with 4 to 9 violations.
Evolution tied with a fixed rule and random search; violations never fell.

\textbf{Violation and success.} The released physical predicates define the
violation: block drops, top-block lateral displacement, boundary exit. Nothing is
shrunk. Task return is sustained safe success, the released predicate held for two
further seconds without violation.

\textbf{Main condition.} The filter runs in shadow mode and records what it would
have done. The adapter acts on the actor's proposal; the executed command is what
the adapter returns, clipped to the legal box. The predicates label outcomes and
veto nothing. The same frozen adapters are rerun with the filter live as the
separate shielded comparison.

\textbf{Adapter.} A bounded residual on the seven coordinates,
$a_t=\operatorname{clip}(a_t^{\mathrm{base}}+B\odot\tanh u_t)$, with bound 0.5 on
the arm and 2.0 on the gripper so a release can become a hold. One decision per
five controls. Inputs: end-effector pose, finger positions, the proposal, three
steps of history and time, 44 in all; hidden 8; 423 parameters. Zero residual
must reproduce native execution bitwise.

\textbf{Arms.} CMA-ES under three ranking rules: task only, penalty with a weight
grid, and feasibility-first at $\epsilon$. Parameter random search as the honesty
control. PPO-Lagrangian as the RL reference with a first-violation cost. Native
actor, untrained adapter, release rule and live filter as references. Eight
episodes per candidate, limit 0.125, population 16, six generations, three
repeats; then the dial at limits 0, 0.125 and 0.25.

\textbf{Panels.} Fresh seed ranges at 3.1, 3.2 and 3.3 million for training,
selection and the 64 reserved final starts, never used before. Policy seed equals
start seed. Overlap with the checkpoint's training data is unknown and declared.

\textbf{Budget.} 44 seconds per episode as measured, half of it simulator start-up.
The full comparison is about 107 lane-hours. Qualification measures two levers,
in-process resets and several environments per process, before anything is
frozen. If neither qualifies, repeats drop to two and PPO moves to a second
allocation.

\textbf{Order.} Qualification: rebuild the platform, reproduce the task 9 native
outcome, prove zero-residual parity, measure adapter authority with generic
genotypes, measure throughput, anchor the references. Freeze. Main comparison.
Focused ablations: trained versus untrained, feasibility versus task-only,
without finger positions. Dial study, separately frozen.

\textbf{What would establish the claim.} Feasibility-first controllers whose
measured risk tracks the requested limit, with completion near the hand-designed
references, while penalty evolution needs a weight sweep to get there and the RL
reference does not get there at matched budgets. A clean negative is also a
result and costs less.

\textbf{Limits.} Migrated simulator, not the original stack or a real robot.
Finger positions are extra information shared by every arm. Equal episode counts
are not equal work; actual controls and queries are reported. Zero observed
violations is not zero risk. This is a working design, not an execution
authorisation.

{\small Detail: \href{BlockPlucking.md}{\texttt{BlockPlucking.md}},
\href{Qualification.md}{\texttt{Qualification.md}}.}
\end{document}
